\section{Event controls}
\label{sec:pfw-event-controls}
\index{PFW!event controls}
\index{events!PFW input}
\providecommand{\eventidx}[1]{\index{MPMEvents!#1@\texttt{#1}}}
\providecommand{\eventattridx}[1]{\index{MPMEvents attribute!#1@\texttt{#1}}}
\pfwidx{useEvents}\pfwidx{mpmEventsString}
\eventattridx{startTime}\eventattridx{endTime}

This section describes how a PFW input file writes solver events.  The solver-side algorithms and required/optional event inputs are documented in Section~\ref{sec:events}, and the generated schema-derived event catalogue is in Appendix~\ref{app:events}.  The PFW control layer is intentionally simple: users enable events with \texttt{pfw["useEvents"] = 1} and place XML event child nodes in \texttt{pfw["mpmEventsString"]}.  PFW writes the surrounding \texttt{<MPMEvents>} block inside the generated \texttt{SolidMechanics\_MPM} solver block.

Use current event attributes, \texttt{startTime} and \texttt{endTime}, in new inputs.  Older PFW examples sometimes include an explicit \texttt{<MPMEvents>} wrapper or use \texttt{time} and \texttt{interval}; the current PFW normalization layer accepts many of those legacy snippets, strips the wrapper, and maps legacy attributes to current event names.  New files should avoid that compatibility path so that the generated GEOS XML is explicit and easy to review.

\begin{lstlisting}[language=Python,caption={General PFW event block pattern.}]
pfw["useEvents"] = 1
pfw["mpmEventsString"] = """
<InitializeStress
  startTime="0.0"
  endTime="1.0e-6"
  targetRegion="all"
  pressure="10.0e6"/>
"""
\end{lstlisting}

Events are evaluated by the solver event manager near the beginning of a time step, before the particle-to-grid transfer sequence, except for the late \texttt{TransformParticles} pass described in Section~\ref{subsec:event-transform-particles}.  Multiple event nodes may be placed in one string.  Event order in the input should follow the intended physical sequence, especially when two events share the same activation time and modify related state.

\subsection{Examples by event type}
\label{sec:pfw-event-examples-by-type}
\index{PFW example!MPM events}

The following fragments show the typical PFW input form for every event type registered in the reviewed source snapshot.  They are intentionally minimal: a production input usually combines these snippets with the boundary controls in Section~\ref{sec:pfw-boundary-controls}, the solver controls in Section~\ref{sec:pfw-solver-controls}, the material assignments in Section~\ref{sec:pfw-materials}, and the output/diagnostic controls in Sections~\ref{sec:pfw-diagnostics} and~\ref{sec:pfw-output-controls}.

\subsubsection{\texttt{Anneal}}
\eventidx{Anneal}\eventattridx{targetRegion}

\texttt{Anneal} acts on a particle region and relaxes the deviatoric part of the stored material stress history while preserving the hydrostatic component.  It is most often used after a consolidation or compaction stage to remove residual shear stress before a subsequent loading branch.

\begin{lstlisting}[language=Python,caption={Anneal event.}]
pfw["useEvents"] = 1
pfw["mpmEventsString"] = """
<Anneal
  startTime="1.0e-3"
  endTime="1.2e-3"
  targetRegion="all"/>
"""
\end{lstlisting}

\subsubsection{\texttt{BodyForceUpdate}}
\eventidx{BodyForceUpdate}\eventattridx{bodyForce}

\texttt{BodyForceUpdate} overwrites the solver-level body-force vector at the event time.  The updated vector is then used by later particle-to-grid body-force projection.

\begin{lstlisting}[language=Python,caption={Body-force update event.}]
pfw["useEvents"] = 1
pfw["mpmEventsString"] = """
<BodyForceUpdate
  startTime="0.0"
  bodyForce="{0.0, -9.81, 0.0}"/>
"""
\end{lstlisting}

\subsubsection{\texttt{BoreholePressure}}
\eventidx{BoreholePressure}\eventattridx{boreholeRadius}\eventattridx{startPressure}\eventattridx{endPressure}\eventattridx{interpType}

\texttt{BoreholePressure} writes a pressure-like background stress in a circular borehole region aligned with the \(z\) axis.  The pressure is ramped between the start and end values using \texttt{interpType}: \texttt{0} for linear, \texttt{1} for cosine, and \texttt{2} for smoothstep interpolation.

\begin{lstlisting}[language=Python,caption={Borehole-pressure event.}]
pfw["useEvents"] = 1
pfw["mpmEventsString"] = """
<BoreholePressure
  startTime="0.0"
  endTime="2.0e-4"
  boreholeRadius="0.015"
  startPressure="5.0e6"
  endPressure="25.0e6"
  interpType="1"/>
"""
\end{lstlisting}

\subsubsection{\texttt{CohesiveZone}}
\eventidx{CohesiveZone}\eventattridx{regionNames}\eventattridx{constitutiveModels}\eventattridx{czTags}\eventattridx{czVolumeNormalization}\eventattridx{computeNormalsAndPositions}\eventattridx{normalsAndPositionsMethod}\eventattridx{czSurfaceDisplacementUpdate}

\texttt{CohesiveZone} creates cohesive-zone regions at run time and associates each region with a cohesive constitutive model and a cohesive-zone tag.  New inputs should use explicit arrays for \texttt{regionNames}, \texttt{constitutiveModels}, and \texttt{czTags}; the arrays must have identical length.  The surface-normal and surface-position options are described in Sections~\ref{sec:cohesive-zone-implementation} and~\ref{sec:contact-surface-detection-spacing}.

\begin{lstlisting}[language=Python,caption={Cohesive-zone insertion event.}]
pfw["useEvents"] = 1
pfw["mpmEventsString"] = """
<CohesiveZone
  startTime="0.0"
  endTime="1.0"
  regionNames="{czInterface}"
  constitutiveModels="{czLaw}"
  czTags="{1}"
  czVolumeNormalization="1"
  computeNormalsAndPositions="1"
  normalsAndPositionsMethod="LogisticRegression"
  czSurfaceDisplacementUpdate="TypeB"/>
"""
\end{lstlisting}

\subsubsection{\texttt{ConfiningPressure}}
\eventidx{ConfiningPressure}\eventattridx{confiningPressureBoxMin}\eventattridx{confiningPressureBoxMax}

\texttt{ConfiningPressure} applies a virtual confining pressure over a rectangular box.  It is commonly paired with \texttt{InitializeStress} so that particles begin with a compatible hydrostatic stress and the boundary/background pressure is then held or ramped over a short start-up interval.

\begin{lstlisting}[language=Python,caption={Confining-pressure event.}]
pfw["useEvents"] = 1
pfw["mpmEventsString"] = """
<InitializeStress
  startTime="0.0"
  endTime="1.0e-6"
  targetRegion="all"
  pressure="10.0e6"/>
<ConfiningPressure
  startTime="0.0"
  endTime="1.0e-3"
  confiningPressureBoxMin="{-0.02, -0.02, -0.001}"
  confiningPressureBoxMax="{ 0.02,  0.02,  0.001}"
  startPressure="10.0e6"
  endPressure="10.0e6"
  interpType="1"/>
"""
\end{lstlisting}

\subsubsection{\texttt{CrystalHeal}}
\eventidx{CrystalHeal}\eventattridx{healType}

\texttt{CrystalHeal} performs a crystal-oriented healing operation on a target particle region.  The \texttt{healType} input is an integer selector used by the solver.  In current source, \texttt{0} is the default branch, while other integer values activate more restrictive or specialized healing choices; users should pair this event with a verification problem before using a new \texttt{healType} in production.

\begin{lstlisting}[language=Python,caption={Crystal-healing event.}]
pfw["useEvents"] = 1
pfw["mpmEventsString"] = """
<CrystalHeal
  startTime="1.20e-3"
  endTime="1.25e-3"
  targetRegion="all"
  healType="0"/>
"""
\end{lstlisting}

\subsubsection{\texttt{DeformationUpdate}}
\eventidx{DeformationUpdate}\eventattridx{prescribedFTable}\eventattridx{stressControl}\index{F-table!event controls}\index{stress control!event controls}

\texttt{DeformationUpdate} switches prescribed-deformation and stress-control modes during a run.  Use it when one stage of a simulation is driven by an \(F\)-table or boundary \(F\)-table and a later stage should change to a stress-control mode, or vice versa.  The actual deformation table and stress-control gains are solver/boundary inputs described in Section~\ref{sec:pfw-boundary-controls}; the event changes which mode is active.

\begin{lstlisting}[language=Python,caption={Deformation-control update event.}]
pfw["useEvents"] = 1
pfw["mpmEventsString"] = """
<DeformationUpdate
  startTime="0.0"
  endTime="1.0e-3"
  prescribedFTable="1"
  stressControl="{0, 0, 0}"/>
<DeformationUpdate
  startTime="1.0e-3"
  endTime="2.0e-3"
  prescribedFTable="1"
  stressControl="{1, 0, 1}"/>
"""
\end{lstlisting}

\subsubsection{\texttt{FrictionCoefficientSwap}}
\eventidx{FrictionCoefficientSwap}\eventattridx{frictionCoefficient}\eventattridx{frictionCoefficientTable}

\texttt{FrictionCoefficientSwap} changes the scalar friction coefficient and/or the group-pair friction table for later contact calculations.  Use a scalar coefficient for a uniform Coulomb coefficient and a table when different contact-group pairs should have different coefficients.  See Section~\ref{subsec:contact-friction-coefficients} for the group-pair lookup convention.

\begin{lstlisting}[language=Python,caption={Friction-coefficient swap event.}]
pfw["useEvents"] = 1
pfw["mpmEventsString"] = """
<FrictionCoefficientSwap
  startTime="5.0e-4"
  frictionCoefficient="0.10"/>
<FrictionCoefficientSwap
  startTime="1.0e-3"
  frictionCoefficientTable="{ {0.00, 0.30}, {0.30, 0.05} }"/>
"""
\end{lstlisting}

\subsubsection{\texttt{Heal}}
\eventidx{Heal}

\texttt{Heal} resets selected damage-related particle state in a target region.  It is the generic healing event, distinct from the polymer- and crystal-specific variants.  It is useful for staged calculations where a preparation step intentionally removes damage before the next loading stage.

\begin{lstlisting}[language=Python,caption={Generic heal event.}]
pfw["useEvents"] = 1
pfw["mpmEventsString"] = """
<Heal
  startTime="1.20e-3"
  endTime="1.25e-3"
  targetRegion="all"/>
"""
\end{lstlisting}

\subsubsection{\texttt{InitializeStress}}
\eventidx{InitializeStress}\eventattridx{pressure}

\texttt{InitializeStress} initializes the stress state of a particle region to a hydrostatic pressure.  This is typically used in geomechanics, periodic-cell preparation, and pre-stressed packing examples before the main load path begins.

\begin{lstlisting}[language=Python,caption={Hydrostatic stress initialization event.}]
pfw["useEvents"] = 1
pfw["mpmEventsString"] = """
<InitializeStress
  startTime="0.0"
  endTime="1.0e-6"
  targetRegion="all"
  pressure="10.0e6"/>
"""
\end{lstlisting}

\subsubsection{\texttt{InsertPeriodicContactSurfaces}}
\eventidx{InsertPeriodicContactSurfaces}\index{periodic contact surfaces}

\texttt{InsertPeriodicContactSurfaces} inserts contact-capable surfaces on periodic boundaries.  It is useful after a densification or healing stage when the subsequent loading branch should treat lateral periodic interfaces as contact surfaces rather than purely topological periodic copies.

\begin{lstlisting}[language=Python,caption={Periodic-contact-surface insertion event.}]
pfw["useEvents"] = 1
pfw["mpmEventsString"] = """
<InsertPeriodicContactSurfaces
  startTime="1.25e-3"/>
"""
\end{lstlisting}

\subsubsection{\texttt{MachineSample}}
\eventidx{MachineSample}\eventattridx{sampleType}\eventattridx{filletRadius}\eventattridx{gaugeLength}\eventattridx{gaugeRadius}\eventattridx{diskRadius}

\texttt{MachineSample} deletes particles to create a test-fixture shape from an initially simpler particle set.  Current options include dogbone-style, Brazilian-disk-style, and cylindrical machining branches; the relevant optional dimensions depend on \texttt{sampleType}.

\begin{lstlisting}[language=Python,caption={Sample-machining events.}]
pfw["useEvents"] = 1
pfw["mpmEventsString"] = """
<MachineSample
  startTime="0.0"
  sampleType="dogbone"
  gaugeLength="0.020"
  gaugeRadius="0.003"
  filletRadius="0.002"/>
<MachineSample
  startTime="0.0"
  sampleType="brazilianDisk"
  diskRadius="0.010"/>
"""
\end{lstlisting}

\subsubsection{\texttt{MaterialSwap}}
\eventidx{MaterialSwap}\eventattridx{sourceRegion}\eventattridx{destinationRegion}

\texttt{MaterialSwap} transfers particles from one particle region to another.  This changes the constitutive model and particle-region membership while retaining the particle state copied by the solver's swap operation.  Define both regions in PFW, even if one begins empty, so the generated XML contains valid source and destination particle regions.

\begin{lstlisting}[language=Python,caption={Material-region swap event.}]
pfw["useEvents"] = 1
pfw["mpmEventsString"] = """
<MaterialSwap
  startTime="2.0e-4"
  sourceRegion="ParticleRegion1"
  destinationRegion="ParticleRegion2"/>
"""
\end{lstlisting}

\subsubsection{\texttt{PolymerHeal}}
\eventidx{PolymerHeal}

\texttt{PolymerHeal} is the polymer-material healing variant.  It acts on the named target region and is usually scheduled after unloading, annealing, or temperature conditioning in polymer verification problems.

\begin{lstlisting}[language=Python,caption={Polymer-healing event.}]
pfw["useEvents"] = 1
pfw["mpmEventsString"] = """
<PolymerHeal
  startTime="1.20e-3"
  endTime="1.25e-3"
  targetRegion="all"/>
"""
\end{lstlisting}

\subsubsection{\texttt{ResetDeformationGradient}}
\eventidx{ResetDeformationGradient}\index{deformation-gradient reset!event controls}

\texttt{ResetDeformationGradient} requests one solver cleanup pass in which the deformation gradient of scaled surface particles is reset as described in Section~\ref{sec:robustness-controls}.  It is useful after geometry or contact-surface operations that should restart the local finite-domain particle shape from a well-conditioned state.

\begin{lstlisting}[language=Python,caption={Reset-deformation-gradient event.}]
pfw["useEvents"] = 1
pfw["mpmEventsString"] = """
<ResetDeformationGradient
  startTime="1.0e-3"/>
"""
\end{lstlisting}

\subsubsection{\texttt{TemperatureProfile}}
\eventidx{TemperatureProfile}\index{temperature table!event controls}

\texttt{TemperatureProfile} copies the solver's domain temperature table value to active particles and cohesive-zone points during the active event window.  The event itself only needs the time window; the temperature history is supplied through the solver/domain temperature-table controls used by the thermal material model.

\begin{lstlisting}[language=Python,caption={Temperature-profile event.}]
pfw["useEvents"] = 1
# The domain temperature table is supplied by the thermal solver controls.
pfw["temperatureTable"] = "{ {0.0, 293.0}, {1.0, 373.0} }"
pfw["mpmEventsString"] = """
<TemperatureProfile
  startTime="0.0"
  endTime="1.0"/>
"""
\end{lstlisting}

\subsubsection{\texttt{TemperatureRamp}}
\eventidx{TemperatureRamp}\eventattridx{startTemperature}\eventattridx{endTemperature}

\texttt{TemperatureRamp} is registered as an event input class and takes a start temperature, end temperature, and interpolation type.  In the reviewed solver snapshot, the event is not dispatched by the active event loop, so it should be treated as a reserved/future event until a regression test demonstrates run-time execution.  Use \texttt{TemperatureProfile} for current temperature-table-driven particle temperatures.

\begin{lstlisting}[language=Python,caption={Registered but currently not dispatched temperature-ramp event.}]
pfw["useEvents"] = 1
pfw["mpmEventsString"] = """
<TemperatureRamp
  startTime="0.0"
  endTime="1.0"
  startTemperature="293.0"
  endTemperature="373.0"
  interpType="2"/>
"""
\end{lstlisting}

\subsubsection{\texttt{TransformParticles}}
\eventidx{TransformParticles}\index{particle transformation!event controls}

\texttt{TransformParticles} is handled by a later event pass after the deformation-gradient update.  It is used by cohesive-zone and finite-domain particle workflows that need to transform stored particle geometry after an initial configuration has been established.

\begin{lstlisting}[language=Python,caption={Transform-particles event.}]
pfw["useEvents"] = 1
pfw["mpmEventsString"] = """
<TransformParticles
  startTime="1.0"/>
"""
\end{lstlisting}

\subsection{Sequential staged event example}
\label{sec:pfw-sequential-events-example}
\index{PFW example!sequential events}
\index{annealing!sequential event example}
\index{healing!sequential event example}
\index{uniaxial stress!event staging}

A staged load path often combines boundary-control tables with event controls.  The following example sketches a common sequence: compress a sample using the active \(F\)-table or moving-boundary controls, anneal the deviatoric stress, heal the damage state, insert lateral periodic contact surfaces, and then unload/reload in tension under uniaxial-stress control.  The event block only schedules state changes.  The actual compression and tensile strain histories must be supplied by the \(F\)-table or boundary controls in Section~\ref{sec:pfw-boundary-controls}.

\begin{lstlisting}[language=Python,caption={Sequential events for compression, annealing, healing, periodic-contact insertion, and tensile loading.}]
# Stage times.
t_compress_end = 1.00e-3
t_anneal_end   = 1.10e-3
t_heal_end     = 1.20e-3
t_periodic     = 1.20e-3
t_tension_end  = 2.50e-3

# Boundary/F-table controls are defined elsewhere in the pfw input.
# For example, the F-table may compress during [0,t_compress_end]
# and then prescribe axial extension during [t_heal_end,t_tension_end].
pfw["prescribedFTable"] = 1
pfw["stressControl"] = [0, 0, 0]

pfw["useEvents"] = 1
pfw["mpmEventsString"] = f"""
<DeformationUpdate
  startTime="0.0"
  endTime="{t_compress_end}"
  prescribedFTable="1"
  stressControl="{{0, 0, 0}}"/>

<Anneal
  startTime="{t_compress_end}"
  endTime="{t_anneal_end}"
  targetRegion="all"/>

<Heal
  startTime="{t_anneal_end}"
  endTime="{t_heal_end}"
  targetRegion="all"/>

<InsertPeriodicContactSurfaces
  startTime="{t_periodic}"/>

<DeformationUpdate
  startTime="{t_heal_end}"
  endTime="{t_tension_end}"
  prescribedFTable="1"
  stressControl="{{1, 0, 1}}"/>
"""
\end{lstlisting}

In this sketch the first \texttt{DeformationUpdate} leaves all three directions prescribed by the deformation table.  The second update keeps the axial deformation history prescribed but turns on stress control in the two lateral directions, producing a uniaxial-stress tensile branch when the boundary/F-table and stress-control parameters are chosen consistently.  For polymer workflows, replace \texttt{Heal} with \texttt{PolymerHeal}; for crystal workflows, replace it with \texttt{CrystalHeal} and supply the desired integer \texttt{healType}.

\subsection{Event-control checklist}
\label{sec:pfw-event-control-checklist}
\index{events!checklist}

Before running a staged event input, check the generated XML and verify the following points:
\begin{itemize}[leftmargin=*]
\item \texttt{pfw["useEvents"] = 1} is set and \texttt{pfw["mpmEventsString"]} is not empty.
\item The string contains event child nodes only.  PFW writes the surrounding \texttt{<MPMEvents>} block.
\item New inputs use \texttt{startTime} and \texttt{endTime}; legacy \texttt{time}/\texttt{interval} snippets have been replaced.
\item Region names such as \texttt{ParticleRegion1}, \texttt{all}, or a custom region name match the generated \texttt{ParticleRegion} names.
\item Cohesive-zone event arrays \texttt{regionNames}, \texttt{constitutiveModels}, and \texttt{czTags} have the same length, and the named cohesive constitutive models are present in the constitutive block.
\item Deformation and stress-control events are consistent with the boundary, \(D\)-table, or \(F\)-table inputs in Section~\ref{sec:pfw-boundary-controls}.
\item Events that modify contact behavior are paired with plot or CSV diagnostics from Sections~\ref{sec:pfw-diagnostics} and~\ref{sec:pfw-output-controls} for at least one short verification run.
\end{itemize}
